\fig{fv:w1:state}{相图以速度为纵坐标，时间增加时按顺时针沿状态圆运动}{\begin{tikzpicture}

\draw[axis](-2.5,0)--(2.8,0)node[right]{$x/A$};\draw[axis](0,-2.5)--(0,2.7)node[above]{$v/(\omega A)$};\draw[blue,thick](0,0)circle(2);
\fill[red](1,1.732051)circle(.05)node[right]{$P_+:v>0$};\fill[red](1,-1.732051)circle(.05)node[right]{$P_-:v<0$};\draw[guide](1,-1.732051)--(1,1.732051);
\draw[flow](1.414214,1.414214)arc(45:10:2);\node[below]at(2,0){$1$};\node[below]at(-2,0){$-1$};\node[left]at(0,2){$1$};\node[left]at(0,-2){$-1$};\node[below left]at(1,0){$1/2$};
\end{tikzpicture}}
图\ref{fv:w1:state}的两轴均已按振幅归一化，圆上的每一点代表一个瞬时状态。$x=A\cos\psi$、$v=-\omega A\sin\psi$给出 $x^2/A^2+v^2/(\omega^2A^2)=1$；右半圆的上、下两点位移相同，速度却反向。给定 $x_0,v_0$后，应同时用余弦和正弦确定初相象限。时间增大时状态沿顺时针转动；它是状态图，不是质点在空间中的圆形轨迹。

参考圆的纵坐标是$A\sin\psi=-v/\omega$，相图的纵坐标却是$v/\omega$，二者反号。因此图\ref{fig:w1:circle}的参考点逆时针转动，图\ref{fv:w1:state}的状态点顺时针转动；同一物理运动没有改变，只是纵坐标换成了相反号的物理量。
